% IOE BE Project LaTeX Template — adapted from Santosh Giri's template
% (as used by Shlok Dai's Mid-Term Defense, located in
% Documentation/97-Educational_Contents/Shlok_Dai_Mid_defense/).
% EdumentX-specific academic content is adapted from the Minor Project
% documentation at Documentation/97-Educational_Contents/EdumentX_Proposal_Minor/.
% Mid-defense chapters (Task Accomplished + Tasks Remaining + Output)
% sit between the methodology chapter and the references.

\documentclass[12pt]{report}
\usepackage{xcolor}
\usepackage[utf8]{inputenc}
\usepackage{natbib}                                % referencing
\usepackage[a4paper,left=1.4in,right=1.0in,top=1in,bottom=1in]{geometry}
\usepackage{float}                                 % [H] figure placement

% Digital-copy hyperref (color links so nav panel is usable):
\usepackage[hidelinks]{hyperref}

\linespread{1.3}                                   % tighter than 1.5; cleaner pages
\usepackage{graphicx}                             % figures
\usepackage{rotating}                             % landscape figures
\usepackage{amsmath}                              % math
\usepackage{amssymb}                              % extra math symbols
\usepackage{titlesec}                             % chapter/section formatting
\usepackage{fontspec}                             % XeLaTeX/LuaLaTeX font selection
% Times New Roman on Windows; Liberation Serif on the TeX Live 2023
% CI toolchain. \IfFontExistsTF picks the available one.
\IfFontExistsTF{Times New Roman}{%
  \setmainfont{Times New Roman}%
}{%
  \setmainfont{Liberation Serif}%
}
\usepackage{enumitem}                             % nicer itemize spacing
\usepackage{longtable}                            % multi-page tables
\usepackage{booktabs}                             % better-looking table rules
\usepackage{microtype}                            % better hbox/line breaking
\usepackage{ragged2e}                             % \RaggedRight to fight overfull hbox

% --- Spacing / line breaking (tighter than Shlok_Dai defaults) ----
\setlength{\parskip}{0.5\baselineskip plus 1pt minus 1pt}
\setlength{\parindent}{0pt}
\setlength{\emergencystretch}{2em}                % keep paragraphs off the right margin
\sloppy                                           % tolerate white-space overflow rather than warn
\tolerance=2000
\hbadness=10000
\hfuzz=2pt                                        % treat <2pt overfull as silent

% --- Chapter / section spacing -------------------------------------
\titlespacing*{\chapter}{0pt}{1.0\baselineskip}{0.8\baselineskip}
\titlespacing*{\section}{0pt}{0.7\baselineskip}{0.3\baselineskip}
\titlespacing*{\subsection}{0pt}{0.5\baselineskip}{0.2\baselineskip}
\titleformat{\chapter}[hang]
  {\normalfont\Large\bfseries\MakeUppercase}
  {\chaptertitlename\ \thechapter.}{1em}{}
\renewcommand{\chaptername}{}

% --- Itemize spacing (so the bulleted Task Accomplished / Tasks
%     Remaining / Output sections don't look gappy) ----------------
\setlist[itemize]{leftmargin=1.2em,itemsep=0.15\baselineskip,topsep=0.2\baselineskip,parsep=0pt}

% --- Image search path ---------------------------------------------
\graphicspath{{Images/}}

% --- safeincludegraphics: thin wrapper that defers to the
%     file path with the Images/ subfolder prefix. graphicx's
%     \includegraphics already consults \graphicspath, so the
%     file lookup Just Works without a fragile \IfFileExists
%     inside a float (which broke [H] placement by returning
%     empty content when a file was missing). ------------------
\providecommand{\safeincludegraphics}[2][]{%
  \includegraphics[#1]{#2}%
}

% --- \outputscreen{file}{caption}{label} : a reusable primitive
%     for the Output chapter. The image lives in the Images/
%     subfolder; graphicx resolves the path via \graphicspath.
\providecommand{\outputscreen}[3]{%
  \begin{figure}[!htbp]
    \centering
    \includegraphics[width=0.85\textwidth]{#1}%
    \caption{#2}\label{#3}%
  \end{figure}%
}

\begin{document}
\begin{titlepage}
    \begin{center}
        \vspace*{1cm}
        {\includegraphics[height=1.5in]{Images/ncelogo.jpg}}\\
        \MakeUppercase{\large{
            {Tribhuvan University}\\
            {Institute of Engineering}\\
            {National College of Engineering}\\
            \vspace{.5cm}
            {A \\Minor Project Proposal\\On\\
            \textbf{EdumentX: Location Based Tutor Discovery Platform}\\
            \vspace{.5cm}
            {{Submitted By:}\\
            Apoorva Paudel (NCE080BCT006)\\
            Asim Poudel (NCE080BCT008)\\
            Dibya Shakya (NCE080BCT013)\\
            Suhan Khadka (NCE080BCT043)}\\
            \vspace{.5cm}
            {{Submitted To:}\\ Department of Electronics \& Computer engineering}\\
        \vspace{1cm}
        Lalitpur, Nepal\\
        February, 2026\\}
        }}
    \end{center}
\end{titlepage}

% --- Front matter --------------------------------------------------
\pagenumbering{roman}
\tableofcontents
\addcontentsline{toc}{chapter}{\numberline{}List of Figures}
\listoffigures
\addcontentsline{toc}{chapter}{\numberline{}List of Tables}
\listoftables
\chapter*{List of Abbreviations}
\addcontentsline{toc}{chapter}{\numberline{}List of Abbreviations}

\begin{tabular}{l l}
AI & Artificial Intelligence \\
API & Application Programming Interface \\
ASCI\-LITE & Australasian Society for Computers in Learning in Tertiary Education \\
BaaS & Backend-as-a-Service \\
CBS & Central Bureau of Statistics (Nepal) \\
DFD & Data Flow Diagram \\
GPS & Global Positioning System \\
GUI & Graphical User Interface \\
HTTP(S) & Hypertext Transfer Protocol (Secure) \\
IDE & Integrated Development Environment \\
IOE & Institute of Engineering \\
JSON & JavaScript Object Notation \\
KNN & K-Nearest Neighbors \\
LBS & Location-Based Service \\
LLM & Large Language Model \\
MFA & Multi-Factor Authentication \\
NCE & National College of Engineering \\
NLP & Natural Language Processing \\
NoSQL & Not only SQL \\
OTP & One-Time Password \\
RAG & Retrieval-Augmented Generation \\
REST & Representational State Transfer \\
SDLC & Software Development Life Cycle \\
SEE & Secondary Education Examination (Grade 10, Nepal) \\
TU & Tribhuvan University \\
UI & User Interface \\
UX & User Experience \\
\end{tabular}

\newpage

% ====================================================================
% CHAPTER 1: INTRODUCTION
% ====================================================================
\chapter{Introduction}
\pagenumbering{arabic}

\section{Background}

Education is one of the most important foundations for personal
growth and social progress. Around the world, families invest
heavily in their children's education, recognizing that formal
schooling alone may not always be enough. When students struggle
with certain subjects or need extra preparation for exams, many
families turn to private tutoring as a way to fill the gaps.
This practice of hiring tutors outside of regular school hours
has become extremely common, as parents look for additional
academic support to help their children perform better.

Private tutoring, commonly referred to as shadow education, has
become a widespread practice as formal schooling systems often
fail to meet the diverse learning needs of students
\citep{Bray2007Shadow}. This form of supplementary education
has grown particularly fast in developing countries, where
public school resources are limited and teacher-student ratios
remain inadequate \citep{Bray2012Shadow}. Research on shadow
education as a global phenomenon confirms that it has expanded
across all continents and income levels, becoming a parallel
curriculum that operates alongside formal schooling
\citep{Kim2019Shadow}.

Nepal's education system has made notable progress over the past
decade. According to the National Population and Housing Census
2021, the overall literacy rate reached 76.3 percent among the
population aged five years and above, up from 65.9 percent in
2011. The country has also experienced rapid urbanization during
this period, with the urban population share jumping from 17.07
percent in 2011 to 66.17 percent in 2021 \citep{CBS2023}.
Despite these improvements, challenges within the mainstream
education system persist. Overcrowded classrooms, inconsistent
teaching quality, and limited resources in public schools
continue to push families toward private tutoring. In cities
like Kathmandu, dual-income households have become increasingly
common, leaving working parents with limited time for direct
academic supervision of their children. As a result, many urban
families look for home tutors as a practical solution.

The phenomenon of shadow education is deeply rooted in Nepal's
educational landscape. Thapa's study on private and public
school competition in Nepal highlights how supplementary
tutoring has become an established practice, with families
across income levels relying on it to fill gaps left by formal
schooling \citep{Thapa2013Private}. Studies further indicate
that approximately 68 percent of Grade 10 students in Nepal
receive some form of private tutoring, confirming how
widespread this practice has become at the secondary level.

Despite this growing demand, Nepal's tutoring market remains
fragmented and informal. Parents and students typically rely on
word-of-mouth recommendations or unstructured social media
groups to find tutors. These methods offer no verification of
tutor credentials, no proximity-based filtering, and no
intelligent matching. With smartphone penetration in Nepal
reaching 65 to 71 percent, the digital infrastructure exists to
address these problems through a dedicated mobile platform.

This mid-term report presents EdumentX, a location-based
mobile application designed to modernize home tutor discovery
in Kathmandu Valley. The platform enables students and parents
to find nearby verified tutors through an interactive map,
enroll in long-term monthly plans, and receive personalized
recommendations through an AI-powered RAG chatbot.

\section{Problem Statement}

The process of finding a home tutor in Nepal is currently slow,
unreliable, and frustrating for everyone involved. Both students
seeking academic help and tutors looking for teaching
opportunities face significant difficulties due to the lack of
any organized system.

The core problems are:

{\renewcommand{\labelenumi}{(\roman{enumi})}
\begin{enumerate}
\item There is no way for parents to check whether a tutor is
actually qualified or trustworthy before hiring them. Most
families rely on recommendations from people they know, which
limits their options and gives no real measure of teaching
quality or safety.

\item Independent tutors, including university students looking
for part-time work, retired teachers, and subject specialists,
have no proper platform to advertise their services. Many end
up going through agencies that take large commissions, while
others simply cannot find students outside their own
neighborhood.

\item The tools people currently use to find tutors, mainly
WhatsApp groups and Facebook pages, are not designed for this
purpose. Posts are unorganized, location details are vague
(just area names with no actual distance), there is no way to
filter or search, and there is no system to verify anyone's
identity or qualifications.
\end{enumerate}
}

At present, no mobile application exists in Nepal that lets
users find tutors on a map, filter by distance, or get
AI-based suggestions. EdumentX aims to solve these problems
by offering proximity-based search on an interactive map, a
verification system where admins review tutor documents and
award Blue Tick badges, a structured monthly enrollment
process, and a chatbot that pulls live tutor data to give
personalized recommendations.

\section{Objectives}

\begin{itemize}
\item To develop a location-based mobile application that enables
students and parents to discover and connect with nearby home
tutors in Kathmandu Valley using map-based proximity search.

\item To integrate an AI-powered RAG chatbot that allows users to
find suitable tutors through natural language queries grounded
in live platform data.

\item To implement a tutor verification system that helps parents
make safe and informed hiring decisions.
\end{itemize}

\section{Scope}

EdumentX is a role-based mobile platform for connecting students
and parents with home tutors based on location, subject, and
monthly enrollment preferences. The scope of this project
includes the following:

\begin{itemize}
\item \textbf{Target Users:} Students, parents, independent home
tutors, and platform administrators.

\item \textbf{Geographic Focus:} Kathmandu Valley, where
smartphone penetration and internet connectivity are
sufficient to support the platform.

\item \textbf{Service Type:} Physical (in-person) home tuition
services only; no online or video-based classes.

\item \textbf{Academic Levels:} School and college students at
SEE, +2, and undergraduate levels.

\item \textbf{Platform:} Cross-platform mobile application for
Android and iOS built with React Native (Expo).

\item \textbf{Discovery Methods:} Map-based search with proximity
ranking and an AI chatbot for natural language queries.

\item \textbf{Enrollment Model:} Long-term monthly enrollment
plans (1, 3, 6, or 12 months) managed through structured
request and response flows.

\item \textbf{Verification:} Blue Tick badge system for tutor
credential verification by admin.

\item \textbf{Group Learning:} Support for group tuition batches
created by tutors for shared, cost-effective learning.
\end{itemize}

% ====================================================================
% CHAPTER 2: LITERATURE REVIEW
% ====================================================================
\chapter{Literature Review}

\section{Related Work}

This section reviews recent academic research on location-based
tutor matching systems, AI-powered educational platforms, and
identifies the gaps that EdumentX aims to address.

A 2022 study presented at ASCILITE introduced MENTOR, an
intelligent mobile peer tutoring application designed for both
face-to-face and remote tutoring \citep{Chung2022MENTOR}. The
system focused on matching peer tutors with students based on
subject expertise and availability. While it demonstrated the
value of mobile platforms for tutoring, it did not include any
location-based features such as map visualization or proximity
ranking. The matching was based solely on academic parameters
without considering geographic distance between users.

A 2026 study published in the Journal of Emerging Technologies
and Innovative Research described a Flutter-based mobile
application for finding local tutors in India
\citep{Sharma2026LocalTutor}. The system used Google Maps API
for proximity-based tutor search and Firebase as the backend
database. Students could filter tutors by subject, availability,
and location. The application calculated distances between
users using the Haversine formula and included features for
booking and rating tutors. However, the researchers noted that
AI-based recommendations were left as future scope. The entire
matching process relied on manual filtering by the user, with
no intelligent ranking or conversational search capability.

In 2025, researchers from Malaysia developed iTeachU, a system
that combined AI with geolocation for tutor matching
\citep{Shari2025iTeachU}. The platform used content-based
filtering along with the K-Nearest Neighbors algorithm,
incorporating Haversine distance calculations to rank tutors
by both preference match and physical proximity. It was built
using Flutter for the frontend and Flask for the backend. While
the system demonstrated effective ranking within the Melaka
region, its AI component was entirely rule-based. Users could
not interact with the system using natural language. Instead,
they had to select from predefined filters and dropdown menus,
making the experience less flexible than a conversational
interface.

All of these studies represent important steps toward
technology-driven tutor discovery, but several gaps remain
when considering the Nepali context. No existing research or
platform provides geospatial visualization of tutors on
interactive maps with service area coverage for Nepal.
City-level dropdown menus, which are used by most platforms,
are not useful in dense urban areas like Kathmandu where a
tutor two kilometers away is very different from one ten
kilometers away. Independent part-time tutors such as
university students and retired teachers lack a proper digital
channel to reach nearby students, as existing systems tend to
favor agency-managed tuition centers. None of the reviewed
systems offer a natural language interface where users can
simply type what they need instead of navigating through
multiple filter options. While some platforms claim AI-based
matching, none combine proximity, qualifications, availability,
and budget into a single intelligent ranking powered by a
large language model. The current methods available in Nepal,
primarily social media groups, require significant manual
effort to browse posts, contact agencies, and compare options
without any automated assistance.

EdumentX brings map-based tutor visualization, a
conversational interface for natural language queries, a
structured verification system with Blue Tick badges, and a
focus on empowering independent tutors to connect with nearby
students efficiently.

\section{Related Theory}

\subsection{Location-Based Services}

Location-Based Services deliver context-aware services using a
device's real-time geographic position obtained through GPS,
cellular triangulation, or Wi-Fi. In EdumentX, location-based
services power the proximity-based tutor discovery and map
visualization features. Tutor coordinates are stored as
Firestore GeoPoint fields, which allow geospatial queries to
be executed directly from the client application.

\subsection{Haversine Distance Formula}

The Haversine formula computes the great-circle distance
between two points on Earth's surface, accounting for the
planet's curvature. EdumentX uses this formula to rank tutors
by their true distance from the student. The computation
proceeds as follows:

\begin{equation}
a = \sin^2\!\left(\frac{\Delta\phi}{2}\right) +
\cos(\phi_1)\cdot\cos(\phi_2)\cdot\sin^2\!\left(\frac{\Delta\lambda}{2}\right)
\end{equation}
\begin{equation}
c = 2\,\arctan2\!\left(\sqrt{a}\,,\,\sqrt{1-a}\right), \qquad
d = R \cdot c
\end{equation}

\noindent where $\phi_1, \phi_2$ are latitudes in radians,
$\lambda_1, \lambda_2$ are longitudes in radians and
$\Delta\phi$ and $\Delta\lambda$ are the coordinate
differences, $R \approx 6{,}371\;\text{km}$, and $d$ is the
output distance. Degree values from GeoPoints are converted
using $\phi_{\text{rad}} = \phi_{\text{deg}} \times
\frac{\pi}{180}$. For the 1 to 20 km range typical in
Kathmandu Valley, this formula provides meter-level accuracy.

\subsection{Geofencing and Radius-Based Filtering}

Before applying the Haversine formula, EdumentX uses a
rectangular bounding box pre-filter to narrow down the
candidate set and reduce computation
\citep{Shevchenko2023Geofencing}. For a given search radius
$r$ and student position $(\phi_s, \lambda_s)$:

\begin{equation}
\Delta\phi_{\max} = \frac{r}{R} \cdot \frac{180}{\pi}, \qquad
\Delta\lambda_{\max} = \frac{r}{R \cdot \cos(\phi_s)} \cdot \frac{180}{\pi}
\end{equation}

\noindent Only tutors whose coordinates fall inside this
bounding box proceed to the exact Haversine distance check,
which keeps query response times low.

\subsection{K-Nearest Neighbors Algorithm}

KNN is a non-parametric algorithm used in EdumentX to rank
tutors by overall suitability. Each tutor is represented as a
feature vector $\mathbf{t} \in \mathbb{R}^5$:

\begin{equation}
\mathbf{t} = [\,d\;, r\;, q\;, e\;, \overline{v}\,]
\end{equation}

\noindent where $d$ is the Haversine distance, $r$ is the
monthly rate, $q$ is the qualification score, $e$ is years of
experience, and $\overline{v}$ is the average rating. All
features are scaled to the range $[0, 1]$ using min-max
normalization: $x_{\text{norm}} = \frac{x - x_{\min}}{x_{\max} -
x_{\min}}$. The similarity between a student query vector
$\mathbf{s}$ and a tutor vector $\mathbf{t}$ is measured by
weighted Euclidean distance:

\begin{equation}
D(\mathbf{s}, \mathbf{t}) =
\sqrt{\sum_{j=1}^{n} w_j \cdot \left(s_j - t_j\right)^2},
\qquad \text{Top-}k = \arg\min_{\mathbf{t} \in \mathcal{T}}\;
D(\mathbf{s}, \mathbf{t})
\end{equation}

\noindent The default value of $k = 5$. Feature weights are
assigned based on their importance for home tuition selection:

\begin{table}[H]
\centering
\caption{KNN Feature Weights for Tutor Recommendation}
\vspace{5pt}
\begin{tabular}{|l|l|l|}
\hline
Feature & Symbol & Weight $w_j$ \\
\hline
Geographic Distance & $d$ & 0.35 \\
Monthly Rate        & $r$ & 0.25 \\
Qualification Score & $q$ & 0.20 \\
Average Rating      & $\overline{v}$ & 0.15 \\
Experience (years)  & $e$ & 0.05 \\
\hline
\end{tabular}
\end{table}

\subsection{Retrieval-Augmented Generation}

Retrieval-Augmented Generation is an approach that combines
information retrieval with text generation to produce responses
grounded in actual data rather than relying solely on a
language model's training data. The original RAG framework was
introduced by Lewis et al., who showed that combining a
pre-trained retriever with a sequence-to-sequence generator
significantly improves accuracy on knowledge intensive tasks
\citep{Lewis2021RAG}. In EdumentX, RAG is used to retrieve
matching tutor profiles from Firestore based on the user's
query and inject them as context into the language model's
prompt, ensuring that recommendations are based on real, live
platform data. A recent survey on RAG chatbots in education
confirms that this approach reduces hallucination and improves
the relevance of generated responses in educational contexts
\citep{Swacha2025RAG}.

\subsection{React Native and Expo}

The EdumentX mobile app is built with React Native, which
lets a single codebase produce native iOS and Android
binaries. Expo sits on top of React Native and provides a
managed toolchain for device features such as location,
camera, and push notifications, along with over-the-air
updates that shorten the release cycle.

\subsection{Firebase and Supabase Storage}

Four services form the EdumentX backend. Authentication
handles email/password and Google sign-in. Firestore is the
main NoSQL store; tutor locations are saved as GeoPoint
fields, and Security Rules control who can read or write
which documents. Supabase Storage (1 GB free, no API key
needed) holds verification documents and demo videos. The
Firebase Admin SDK is used out-of-band from a local Node.js
script to grant the administrator role to a fixed list of
emails, since the \texttt{admins} collection is locked to
Admin-SDK writes only.

The project does not rely on Cloud Functions. Distance
ranking, the KNN scoring, and the retrieval step of the
chatbot all run on the client. This keeps the build within
the free Spark tier and avoids the paid Blaze plan that the
team cannot sign up for.

\subsection{OpenStreetMap and Open Geocoders}

To avoid the paid Google Maps Platform API, EdumentX uses
OpenStreetMap raster tiles for map rendering and the
keyless Nominatim service for forward geocoding (turning a
typed city name into latitude and longitude). Distance and
KNN matching are computed on the device, which is fast
enough at the Kathmandu-Valley scale where the target
database size is in the low thousands of tutors.

% ====================================================================
% CHAPTER 3: METHODOLOGY
% ====================================================================
\chapter{Methodology}

\section{System Design}

\subsection{System Architecture}

Figure~\ref{fig:sysarc} shows the overall system architecture of
EdumentX. The architecture has three layers: the client
side, the serverless backend, and external services.

\begin{figure}[H]
\centering
\safeincludegraphics[width=\textwidth]{sysarc.png}
\caption{System Architecture of EdumentX}
\label{fig:sysarc}
\end{figure}

The client side is the React Native mobile application
built with Expo. It serves as the single entry point for all
users and communicates with the backend using the native
Firebase SDK.

The serverless backend runs entirely on Firebase with three
services. Firebase Authentication handles user registration
and email/password plus Google sign-in. Cloud Firestore
stores all structured data including user profiles, tutor
profiles, enrollment requests, and reviews, with tutor
locations stored as GeoPoint fields and protected by
Security Rules. Supabase Storage (1 GB free tier) holds
tutor verification documents and demo videos.

The external services layer includes OpenStreetMap tiles
for map rendering, Nominatim for keyless geocoding, and a
free-tier LLM endpoint for the chatbot. EdumentX does not
depend on Cloud Functions, Google Maps SDK, or any paid
LLM; every external call is either keyless or fits
within a free quota.

\subsection{Use Case Diagram}

Figure~\ref{fig:usecase} shows the use case diagram with three
actors: Students/Parents, Tutors, and Admin.

\begin{figure}[H]
\centering
\safeincludegraphics[width=0.75\textwidth]{usecase.jpeg}
\caption{Use Case Diagram}
\label{fig:usecase}
\end{figure}

Students can register, sign in, search for tutors using
filters or the AI chatbot, view tutor profiles, send
enrollment requests, cancel an active enrollment, and rate
tutors after completing a plan.

Tutors can register, set their availability, upload
documents for verification, manage their profile, respond to
enrollment requests, and create group batches for shared
tuition.

Admin can log in to the admin panel, verify or reject tutor
documents, manage user accounts, and grant Blue Tick badges
to verified tutors.

\subsection{Data Flow Diagram}

Figure~\ref{fig:dfd0} shows the Level 0 Data Flow Diagram.
Users interact with the EdumentX system by sending requests
and receiving responses. The system processes these requests
and stores or retrieves data from Firebase Firestore.

\begin{figure}[H]
\centering
\safeincludegraphics[width=2.8in]{Level_0_DFD.png}
\caption{Level 0 Data Flow Diagram}
\label{fig:dfd0}
\end{figure}

Figure~\ref{fig:dfd1} shows the Level 1 Data Flow Diagram
with detailed system processes including authentication,
tutor search, AI chatbot recommendation, enrollment
management, tutor verification, and rating handling.

\begin{figure}[H]
\centering
\safeincludegraphics[width=\textwidth]{Level_1_DFD.png}
\caption{Level 1 Data Flow Diagram}
\label{fig:dfd1}
\end{figure}

\subsection{RAG Chatbot Workflow}

Figure~\ref{fig:ragworkflow} illustrates the RAG workflow
used by the EdumentX AI chatbot.

\begin{figure}[H]
\centering
\safeincludegraphics[width=\textwidth]{RAG.png}
\caption{RAG-Based AI Chatbot Workflow}
\label{fig:ragworkflow}
\end{figure}

The chatbot works through a six-step pipeline. First, the
student types a natural language query like ``I need a science
tutor near Lalitpur for 3 days a week under Rs 5,000.'' The
full conversation history is included so the chatbot can
handle follow-up questions. Second, the client extracts
structured parameters from the query: subject, location,
budget, days per week, and plan duration. The free-tier
LLM endpoint handles this extraction. Third, the client queries the Firestore
\texttt{tutors} collection, applies the bounding-box
pre-filter followed by exact Haversine distance, and applies
budget and subject filters. Fourth, the top ten matching
tutor profiles are formatted and injected into the LLM's
prompt as context data. Fifth, the language model generates
a response recommending three to five tutors with name, rate,
distance, rating, and verification status. Sixth, the
response is rendered as a chat message in the app. Identical
queries are cached for 15 minutes to keep the free-tier
token budget healthy.

% ====================================================================
% CHAPTER 4: TASK ACCOMPLISHED
% ====================================================================
\chapter{Task Accomplished}

EdumentX is a mobile-first tutor discovery platform for
Kathmandu Valley that helps students and parents find nearby
verified home tutors through an interactive map, an
AI-powered chatbot, and a structured monthly enrollment
process. The following tasks have been accomplished till
now as part of the EdumentX mid-term development.

\section{Project Setup and Documentation:}

A solid planning base was set up before any code was written.
A feature inventory and folder structure for documentation,
reference material, and archived work were created so that
the team could move in a single direction. A team change
log and a project status document are maintained alongside
the repository.

\section{Design System:}

The application's visual design is built around a small set
of design tokens for colors, surfaces, borders, and text
that are registered in the build configuration. Shared UI
primitives such as the button, card, screen, and header
have been extracted for reuse. The styling layer was
migrated off the previous Tamagui setup to plain NativeWind
class names, which removed a large dependency surface and
made the codebase easier to maintain.

\section{User Sign-in and Sign-up:}

EdumentX supports two free sign-in methods. Email and
password registration and login are wired through native
Firebase Authentication, and Google sign-in is wired through
the native Google Sign-In SDK. After an email sign-up, the
user is shown a ``check your inbox'' panel that stays
visible until the email is verified. An ``I've verified''
button on that panel refreshes the session so a user who
returns from the inbox is not stuck on the pending screen.
A routing guard in the root layout sends unverified users
back to the sign-up screen, users without a role to the
role selection screen, and verified users onward to their
dashboard. A bug where an existing user was being routed
back to role selection on every sign-in was identified and
fixed.

\section{Onboarding and Profile:}

After sign-in, a new user lands on a role selection screen
that writes the chosen role to the user record and routes
the user to the matching onboarding form. The student
onboarding form captures name, phone, city, grade, and
subjects. The tutor onboarding form captures name, phone,
city, headline, bio, monthly rate, subjects, and years of
experience. A minimum length check on the city field
prevents the form from being submitted with a one-character
placeholder. Profile photos are uploaded to object storage
and the resulting URL is persisted to the user record so
that the photo survives an app restart.

\section{Dashboards:}

The student home and tutor home screens are wired to live
reads from the database. Each shows the real name and city
of the signed-in user instead of hard-coded placeholder
text. The tutor home additionally shows the verification
status that the admin has assigned.

\section{Profile Tabs:}

The student and tutor profile tabs share the same layout,
which consists of an identity card, editable fields, menu
rows, and a destructive log-out button. The profile tab
allows the user to pick a new profile photo from the device
library. A previously missing log-out affordance on the
tutor profile tab has been added.

\section{Admin Panel:}

The admin panel exposes a home view, a statistics screen, a
verification queue, a user management screen, and an admin
profile setup screen. Because the admin role is sensitive,
it is not assigned from inside the app. A small local admin
script is used to grant the role to a fixed list of emails
through the Firebase Admin SDK, which keeps the underlying
\texttt{admins} collection locked to out-of-band writes
only.

\section{Onboarding Visuals:}

The first two onboarding slides, which introduce the
platform's discovery map and the AI chatbot, have been
re-rendered against the new design palette so that they
match the rest of the application.

\section{Verification:}

The mid-term build passes a clean TypeScript type check and
a clean ESLint run with zero errors. A manual end-to-end
test on an Android device was carried out for the student,
tutor, and admin roles.

% ====================================================================
% CHAPTER 5: TASKS REMAINING
% ====================================================================
\chapter{Tasks Remaining}
\label{chap:remaining}

The following are the tasks which need to be accomplished
between the mid-term defense and the final defense.

\begin{itemize}
\item Tutor Search and Map Discovery
\item AI Chatbot with Live Tutor Data
\item Tutor Verification and Blue Tick Badge
\item Enrollment Request and Approval Flow
\item Ratings and Reviews
\item Admin Dashboard with Real Platform Data
\item Final End-to-End Testing and Demo Build
\end{itemize}

% ====================================================================
% CHAPTER 6: OUTPUT
% ====================================================================
\chapter{Output}

This chapter shows the screens built so far. The current
build is shown in Figure~\ref{fig:mockup1}.

\outputscreen{mockup_1.png}{EdumentX Current Build Mockup}{fig:mockup1}

% ====================================================================
% REFERENCES
% ====================================================================
\renewcommand{\bibname}{References}
\phantomsection
\addtocontents{toc}{\protect\contentsline {chapter}{References}{}{}}
\bibliographystyle{unsrt}
\bibliography{ref}
\end{document}
